\documentclass[10pt]{article}
\setlength\parindent{10pt}
\usepackage{amsmath, amssymb}
\usepackage[margin=1in]{geometry}
\usepackage{mathrsfs}
\usepackage{natbib}
\pagestyle{empty}

\title{Assessing journal and institution influence through reputation recirculation via citations}
\date{Technical Report, \today}
\author{Domingo Docampo and Lawrence Cram}
\begin{document}
\maketitle

\begin{abstract}

Traditional citation-based metrics, like the Journal Impact Factor, measure popularity but are vulnerable to manipulation and fail to account for the prestige of the citing source. Eigenvector-based methods (also called spectral ranking) address the desirability of accounting for the influence of citing journals although they remain within the reflexive journal-to-journal citation system and can be manipulated by orchestrating citations.

 We introduce a novel method, the Journal-Institution Reputation Circulation (JIRC) algorithm, that recirculates reputation between journals and institutions via a two-mode citation network to derive more robust influence metrics.
We construct two normalized citation matrices: one from institutions to journals (Matrix A) and the other from journals to institutions (Matrix B). The algorithm solves for the dominant eigenvector of the combined system \(v = P_B^T P_A^T v\), where \(v\) is the institutional influence vector, and derives the journal influence vector \(w\) as \(w = P_A^T v\). We apply JIRC to a large dataset from the ESI Field of Economic \& Business.

The resulting journal and institution classifications show strong construct validity, aligning with established reputations. Moreover, when used to screen a sample of 330 researchers—equally dividing them into a ``Control'' group (Nobel laureates and prestigious award winners) and a ``Test'' group (authors with highly cited papers but no major awards)—JIRC effectively disentangles influence from notoriety. The top tier of researchers was predominantly from the Control group, while the bottom tier was overwhelmingly composed of members from the Test group. The JIRC algorithm thus provides a robust, manipulation-resistant tool for assessing influence at the journal, institutional, and individual researcher level, offering a significant advancement over purely citation-count-based methods.

{\bfseries Keywords}: Scientometrics, Research Evaluation, Citation Analysis, Reputation, Eigenvector Centrality, Pinski-Narin, PageRank, Journal Influence, Institutional Ranking
 \end{abstract}
---

\section{Introduction}

The review, evaluation and assessment of a researcher's work is a high-stakes activity that shapes professional success in publication, recognition, grant funding, appointment, advancement and financial rewards. The temptation to boost, curate, misrepresent, orchestrate or fake research is accordingly high, countered by regulations, moral/ethical values, and the opprobrium attached to violation of norms \citep{mertonSocialStructureScience1996}. The interplay of temptation and morality leads to the publication of works designed primarily to inflate performance indicators, rather than report findings that advance research. The proportion of such `aberrations' is not well characterised, but the informed estimate by \cite{jamesheathersApproximately172025} that 1/7 of the literature is `fake' corresponds currently to around 700,000 works per year. Efficient and accurate methods to screen the research literature for aberrations at this scale would support well-focused, systematic, organised scepticism that is so important to scientific progress \cite[Chapt. 9]{zimanRealScienceWhat2001}.

Some aberrations in scientific publications are revealed by material within the document, including fake papers \citep{elseFightFakepaperFactories2021}, plagiarism \citep{longRespondingPossiblePlagiarism2009}, data fabrication \citep{brownIssuesDataAnalyses2018}, falsification of methods \citep{simonsohnJustPostIt2013} and image manipulation \citep{bikPrevalenceInappropriateImage2016}. For others, potential aberrations are revealed by bibliometric and reference metadata, including questionable author-editor relations \citep{banksElsevierChallengedJournal2009}, guest authorship \citep{khezrVexingPersistentProblem2022}, coerced references \citep{straubJournalSelfCitationVI2009}, hyper-production \citep{ioannidisFeaturesSignalsPrecocious2024a}, citation cartels \citep{safonScreeningArticlesCitation2025}, and orchestration \citep{evdaimonMetricsDetectSmallScale2024}.

Bibliometric studies are not normally designed to detect aberrations but rather to identify and rate worthy research or reveal the practice of science. Nevertheless, some are applicable to large scale screening for aberrations if one moves from an appreciative to investigatory perspective: e.g. scrutinising high publication or citation counts rather than offering accolades. As discussed by \cite{tahamtanSocialSystemsCitation2022} in their proposal for the Social Systems Citation Theory (SSCT), a bibliometric approach does not need to deal with humans and their motives but rather may begin with and focus on the communication network of research. Once identified, aberrant works may be set aside \citep[e.g. in systematic reviews,][]{munnShouldIncludeStudies2021}, perused through the lens of forensic bibliometrics \citep{mcintoshForensicScientometricsAnEmerging2024}, or subject to deeper but time-consuming textual and peer review \citep{kilicogluBiomedicalTextMining2018}. Aberrant work may also help select subjects for ethnographic or psychosocial studies of researchers' individual motives and collective behaviours \citep{martinsonScientistsBehavingBadly2005}.

Effective large-scale screening for anomalies using journal-level bibliometric indicators has been demonstrated by \cite{kojakuDetectingAnomalousCitation2021}. Their CIDRE (Citation Donors and Recipients) algorithm detects anomalous citation patterns between journals, revealing candidates for possible exclusion from journal indexing services. \cite{chakrabortyIdentificationAnalysisCitation2020} address a similar problem using sets of features to identify anomalies in computer science journals. \cite{evdaimonMetricsDetectSmallScale2024} exhibit author-level screening using the referencing information of tens of millions of articles to quantify the extreme values of field-specific indicator distribution functions that can isolate specific patterns of referencing orchestration. \cite{ioannidisFeaturesSignalsPrecocious2024a},  \cite{ioannidisEvolvingPatternsExtreme2024a}, and \cite{ioannidisQuantitativeResearchAssessment2023a} identify many indicators that potentially herald large-scale manipulation. \cite{ibrahimCitationManipulationCitation2025} exhibit an author's citation concentration index, $c^2$, the largest number $n$ of papers that cite the author at least $n$ times, showing that it can be used as a warning about manipulation and also revealing how pre-print servers may be exploited.

These and similar studies demonstrate that the tails of some frequency distributions (FDs) derived from populations of aggregated citations or references include a higher proportion of aberrant publications. An alternative to examining extrema of FDs is to screen a reference or citation before aggregation, for example by applying a weight.  An article's field-weighted citation index (FWCI) and relative citation ratio (RCR) are familiar indicators that weight citations to adjust for benign referencing variations between fields of research. The reputable citation (RC) method of \cite{safonScreeningArticlesCitation2025} uses an exogenous rating of citing author's institutions to scale citations prior to aggregation at the journal or author level. Applied to the field of Mathematics, a large gap between aggregations of scaled and unscaled citations delineates authors whose elevated citation profile may have little to do with the modest impact and quality of their work.

While the exogenous institution ranking used in the RC method helps to insulate the indicator from anomalous publications, it is challenging to find the large-scale, independent, field-specific institution rankings that the method uses. Accordingly in this paper we use spectral ranking over a multi-partite citation network to establish indicators of journal and institution prestige. These are then conferred on citations by works hosted by the respective sources and institutions. The raw and scaled citation counts are then used to compute an author esteem indicator. In this way we exhibit a screen that promises to reveal low-esteem authors who contribute anomalous works to the research literature.

The paper is organised as follows. First, we introduce working distinctions between influence, status, prestige and esteem, linking influence and prestige to spectral ranking over normalised citation matrices. We present a recirculation algorithm for computing prestige indicators as the spectral rank over a multi-partite citation network projected to journals and institutions. Using the journal and institution prestige indicators we explain how to compute authors' relative esteem. We demonstrate the approach using the field of Economics, Finance and Business, and discuss our findings about the use of this approach as a screen. We conclude with a discussion of the limitations and prospects of this family of methods. 

\section{Prestige, esteem and spectral ranking}
Researcher publication rates (articles/year) vary by orders of magnitude due to many factors and are only weakly linked to the status of their oeuvres \citep{ScientificOutputRecognition, ioannidisEvolvingPatternsExtreme2024a}. Publication citation rates (cites/article/year) also vary by orders of magnitude due to many factors \citep{tahamtanFactorsAffectingNumber2016, tahamtanWhatCitationCounts2019}, including an aspect of status called `influence' or `popularity' \citep{bollenJournalStatus2006a, garfieldeugeneCitationIndexingIts1979}. Similar observations can be made about the citation rates of journals and institutions accumulated through their constituent publications.

That \textit{status} is more than popularity has long been recognised in sociology \citep[e.g.][]{katzNewStatusIndex1953, podolny05}. It will be useful in this work to have terms for two kinds of status indicator: \textit{prestige} for journals and institutions, and \textit{esteem} for researchers. \citep[This distinction reflects discussions in, e.g., []{davisPrinciplesStratification1945, goodeCelebrationHeroesPrestige2022}. \textit{Prestige} is a long-term, relatively stable status indicator that is largest for journals like \textit{Nature} \citep{baldwinMakingNatureHistory2019} and institutions like those at the top of global university rankings \citep{aguilloComparingUniversityRankings2010}. \textit{Esteem} is a short-term, fluid status indicator that could be perceived to rise if a senior researcher becomes tenured at a high-prestige university, and decline if they are censured for questionable research practices \citep{aronsonGainLossEsteem1965, hassanibesheliGainLossEsteem2017}. 
    

According to our usage, the long-standing ISI journal impact factor (JIF) introduced in \cite{garfield1955, garfieldHistoryMeaningJournal2006} is an indicator of \textit{journal prestige}, especially within a discipline. JIF reports the average number of citations per article in a journal, counting citations with equal value and thus ignoring the nature of citing works or journals. \cite{garfieldHistoryMeaningJournal2006} explains that the JIF helps identify the core journals in a field, some by larger size and others by higher citation frequency (higher impact), and suggests that some authors use the JIF to help select an outlet their work since `journals with high impact factors include the most prestigious' (p. 92). Note that \cite{bollenJournalStatus2006a} argue that the JIF is an indicator of popularity not prestige, since it takes no account of the standing of referencing authors. 
 
Many researchers have advocated alternatives to the JIF that take account of the context of a reference. \cite{waltmanRelationEigenfactorAudience2010} categorise alternatives as (a) normalising the cited-side, (b) normalising the citer-side, (c) using an indicator other than citations/article, or (d) recursively weighting citations. For example, the citer-side adjusted Audience Factor \citep{zitt2008} adjusts the JIF for inter-journal differences by boosting or diminishing the value of a journal's references in proportion to the average number of references per article before summing the weighted citations. This compensates for disciplinary differences in the length of reference lists and in the delay before work is cited, both factors that influence the JIF. \cite{habibzadehJournalWeightedImpact2008} propose citer/citing-side scaling that weights a reference by the quotient of the JIFs of citing and cited journal, rescaling this weight to the range 0.1 to 10 over a logit function of the JIF ratio. 

The JIF and many variants focus initially on individual journals characterised by the number, references and citations of their constituent papers. Comparisons to establish the relative standing of journals are made after the JIFs are computed. An alternative view is that journals inhabit a network or social system where relative standing is determined recursively by the whole community. An early study based on this view \citep{narinEVALUATIVEBIBLIOMETRICSUSE1976, pinski, narinStructureBiomedicalLiterature1976} [hereinafter PN] measures the influence of a publishing unit (journal, department, institution, nation) as the stable distribution of `influence weight' (a kind of prestige) that is exchanged through reference lists. \cite{vignaSpectralRanking2016} adopts the generic term \textit{spectral ranking} in a helpful summary of the antecedents and followers of PN.

PN provide a general approach to using citation matrices to spectrally rank \textit{units} such as authors, journals, institutions, or countries that they apply to journals. Here, we expand their approach to consider parallel spectral ranking of two units, namely journals and institutions.

The entities required to build this model comprise sets of $N$ works, $M$ sources, $P$ authors and $Q$ institutions, selected to relate to a specified discipli. Works published in the \textit{census} time interval $T^C$ have a \textit{reference list} pointing to referenced works published in an earlier or overlapping \textit{target} time interval $T^T$ that establishes a binary directed work-work network. Works metadata then establishes adjacency matrices: 

\begin{itemize}
	\item[] \text{Reference Matrix} ($\mathbf{R} \in \{0,1\}^{N \times N}$):
	$R_{ij} = 1$ if work $i$ references work $j$.
	
	\item[] \text{Publishing Matrix} ($\mathbf{B} \in \{0,1\}^{N \times M}$):
	$B_{ws} = 1$ if work $w$ is published in source $s$.
	
	\item[] \text{Authorship Matrix} ($\mathbf{U} \in \{0,1\}^{N \times P}$):
	$U_{wa} = 1$ if work $w$ is written by author $a$.
	
	\item[] \text{Affiliation Tensor} ($\mathcal{V} \in \{0,1\}^{N \times P \times Q}$): 
	$V_{wai} = 1$ if author $a$ lists institution $i$ on work $w$.
	
	\item[] \text{Institutional Presence Matrix} ($\mathbf{Z} \in \{0,1\}^{N \times Q}$):
%	This derived matrix indicates if an institution is present on a work (regardless of which author listed it). \\
	$Z_{wi} = 1$ if $\sum_{a=1}^{P} V_{wai} \geq 1$, else $0$.
\end{itemize}
To map works to institutions, we calculate the work-institution weight matrix $\mathbf{H} \in \mathbb{R}^{N \times Q}$. Depending on the chosen bibliometric counting method, $\mathbf{H}$ is defined in one of two ways: \\
\\
Author Fractionation (Credit institutions based on the share of authors they contribute.)
\begin{align}
	H_{wi} = \sum_{a=1}^{P} \underbrace{\left( \frac{U_{wa}}{\sum_{p=1}^{P} U_{wp}} \right)}_{\text{Author Share}} \underbrace{\left( \frac{V_{wai}}{\sum_{q=1}^{Q} V_{waq}} \right)}_{\text{Affiliation Share}}
\end{align}\
Institution Fractionation (Credit institutions equally, independent of the number of authors.)
\begin{align}
	H_{wi} = \frac{Z_{wi}}{\sum_{q=1}^{Q} Z_{wq}}
\end{align}
Reference flow in this model is described by the block matrix $\mathbf{K} \in \mathbb{R}^{(M+Q) \times (M+Q)}$
\begin{align}
	\mathbf{K} = 
	\begin{pmatrix}
		\mathbf{K}_{SS} & \mathbf{K}_{SI} \\
		\mathbf{K}_{IS} & \mathbf{K}_{II} 
	\end{pmatrix} = 
	\begin{pmatrix}
		\mathbf{B}^\top \mathbf{R} \mathbf{B} & \mathbf{B}^\top \mathbf{R} \mathbf{H} \\
		\mathbf{H}^\top \mathbf{R} \mathbf{B} & \mathbf{H}^\top \mathbf{R} \mathbf{H} 
	\end{pmatrix}
\end{align}
The blocks represent projected reference flows, $\mathbf{K}_{SS}$: Source referencing Source, $\mathbf{K}_{SI}$: Source referencing Institution, $\mathbf{K}_{IS}$: Institution referencing Source, and $\mathbf{K}_{II}$: Institution referencing Institution. 

$\mathbf{K}_{SS}$ is the source citation matrix $\mathbf{C}$ defined by PN. The element $K_{SS} (i,j) = c(i,j)$ is the number of references from journal $i$ to journal $j$ and the number of citations received by $j$ from $i$. Since journals publish different number of articles per year, and articles have reference lists of different lengths, PN scale the exchange from journal $i$ to $j$ as the input-output ratio $c_{i,j}/s_j$, where the numerator is the reference flow from $i$ into $j$ in $T_C$ and the denominator $s_j = \sum_k c_{k,j}$ is the count of all references to $j$ in $T_C$. The final \textit{influence per reference} of each journal, $w_i$, is then determined recursively as the influence-weighted sum of the normalised ratios, $w_i = \sum_k w_k c_{i, k}/s_k$, with the scaling of $w$ set by the condition that the average size-weighted influence is unity, $\sum_i s_i w_i = \sum_i s_i$. The influence per publication is $v_i = \sum_kv_k (c_{i,k}/s_k) (a_k/a_i)$ where $a_i$ is the number of articles published in $T_c$ by journal $i$ \citep[see also][]{palacios-huertaMeasurementIntellectualInfluence2004}.  The Index of New Knowledge (INK) proposed by \cite{paulyReinterpretation145influenceWeight1462008} is the initial, non-recursed stage of the Pinski-Narin algorithm.

\cite{palacios-huertaMeasurementIntellectualInfluence2004} [PHV] adopt an axiomatic approach to the unit ranking problem (their work considers only journals as the unit, but this is not required). Defining reference intensity as the mean number of references per article in a unit, PHV show that the PN algorithm for the influence \textit{per publication} ($v_i$ above) uniquely satisfies four axioms: (1) the influence ratio is not affected by the reference intensity, (2) for units of equal size and reference intensity the influence ratio is equal to $C_{i,j}/C_{j,i}$, (3) in a set of units having the same reference intensity, the removal of a unit does not change the influence ratio, and (4) splitting a unit into subsets does not change the influence ratio. \cite{serranoMeasurementIntellectualInfluence2004} offers an impossibly result to counter PVH's optimism, showing that the inclusion of a large unit with low citation flow to and from the original unit set markedly shifts $v_i$. \cite{waltmanRelationEigenfactorAudience2010} draw the same conclusion by noting that the PN method is not invariant to disciplinary differences, a point also emphasised by PN.

PN used the power iteration to determine the principal eigenvector of the input-output matrix, and applied the method to rank separately Physics and Biology Journals. Chapter 10 of \cite{narinEVALUATIVEBIBLIOMETRICSUSE1976} also presents an application of the method to institutional influence ranking, but does not implement spectral ranking of institutions as units. Rather, the procedure first constructs journal sets and their influence weights (11 fields) and then derives the total influence of around 135 institutions (broadly and by field) by summing the institutional publication count in each journal weighted by the respective influence weight per publication. The algorithmic institutional influences compare favourably with rankings derived from surveys.

Motivated by this work, \cite{gellerCitationInfluenceMethodology1978} considered an alternative exchange of influence or prestige through the ratio of the number of references from $i$ to $j$ over the total number of references from $i$, $\sum_k c_{ik}$. Geller's normalisation yields a row-stochastic matrix, the transition matrix of a Markov chain. The elements of its principal eigenvector $\pi^*_i$ are the long-run probabilities that unit $i$ will be cited. The connection between the PN influence per reference and Geller's probability is $w_i = \pi^*_i \sum_k s_k / s_i$,  the expected number of citations that $i$ will receive, per reference that it gives out. 

\cite{bollenJournalStatus2006a} propose a weighted PageRank algorithm that computes the probability $\pi^*$ via the row-normalised citation matrix, following \cite{brinAnatomyLargescaleHypertextual1998a} by introducing an attenuation factor $0 \le \alpha \le 1$ and a minimum amount of influence (prestige) for each unit ($1/N$ where $N$ is the total number of references).
\begin{align}
	\pi^* = (1-\alpha)/N + \alpha \sum_j \pi^*_k \frac{c_{ij}}{\sum_k c_{ik}}
\end{align}
Clarivate's Eigenfactor journal score \citep{bergstromEigenfactorMeasuringValue2007, westEigenfactorMetricsNetwork2010} and the Scimgo-Scopus journal ranking \citep{gonzalez-pereiraNewApproachMetric2010,guerrero-boteFurtherStepForward2012} also use the row-normalised form, with variants of $c_{ij}$ that exclude journal self-citations and/or adopt different time windows. \cite{vigna} and \cite{waltmanRelationEigenfactorAudience2010} survey connections between historical antecedents of recursive algorithms, and reveal the important connection between recursions and path summations.

\subsection{Prestige for journals and institutions (JIRC algorithm)}

The PN influence per reference $w_i$ is a measure of the influence of a reference (or equivalently a `citation') from $i$ relative to the average, and the ratio $w_i/w_j$ is a measure of the relative influence of a citation from unit $i$ and from unit $j$. Differences between unit size and reference density have been removed. Accordingly, we will herein call $w_i$ the \textit{prestige} of unit $i$ and use its value to represent the \textit{esteem} that a work acquires from a citation by $i$. 

We extend the PN solution to the two-unit source-institution problem using the block matrix $\mathbf{K}$ defined above, with diagonal blocks $K_{SS}$ and $K_{II}$ set to zero. The removal of these blocks excises direct journal-journal and institution-institution reference flows, including all unit self-references. This approach tends to strengthen the `independence' of citing-side contributions to the prestige indicators, since direct journal-journal and institution-institution citation is excluded. 

Introduce the diagonal matrix $\mathbf{D} = \text{diag}(1/s_j)$ whose elements are the reciprocals of the column sums of $\mathbf{K}$, $s_j = \sum_k K_{kj}$, and the PN block matrix 
\begin{align}
\bf{\Gamma} = \mathbf{K}\mathbf{D} = 
	\begin{pmatrix}
	\mathbf{0} & \bf{\Gamma}_{SI} \\
	\bf{\Gamma}_{IS} & \mathbf{0} 
\end{pmatrix}
\end{align}
Provided that $\mathbf{\Gamma_{SI}}$ is a primitive matrix, the prestige values are the components of the principal eigenvector
\begin{align}
\bf{\Gamma}^T \mathbf{w} = \lambda \mathbf{w}
\end{align}
with eigenvalue $\lambda = 1$ and the normalisation 
\begin{align}
\mathbf{w} = \text{N} \mathbf{D} \mathbf{w}  
\end{align}
where $N$ is the total number of references, $N = \sum_k s_k$.

If we partition the principal eigenvector into components corresponding to sources ($\mathbf{w}_S$) and institutions ($\mathbf{w}_I$):
\begin{align}
	\mathbf{w} = 
	\begin{pmatrix}
		\mathbf{w}_S \\
		\mathbf{w}_I
	\end{pmatrix}
\end{align}
the components satisfy the \textit{Journal-Institution Reputation Circulation} (JIRC) system that connects the prestige vectors 
\begin{align}
	\mathbf{\Gamma}_{IS}^\top \mathbf{w}_I &= \mathbf{w}_S \\
	\mathbf{\Gamma}_{SI}^\top \mathbf{w}_S &= \mathbf{w}_I
\end{align}
and either one can be found from
\begin{align}
	(\mathbf{\Gamma}_{IS}^\top \mathbf{\Gamma}_{SI}^\top) \mathbf{w}_S &= \mathbf{w}_S \\
	(\mathbf{\Gamma}_{SI}^\top \mathbf{\Gamma}_{IS}^\top) \mathbf{w}_I &= \mathbf{w}_I
\end{align}
Our JIRC uses the power iteration method
\begin{align}
\mathbf{w}_S^{(k)} = \mathbf{\Gamma}_{IS}^\top \mathbf{\Gamma}_{SI}^\top \mathbf{w}_S^{(k-1)}
\end{align}
yielding the source prestige vector $\mathbf{w}_S$ and by (10) the institution prestige vector $\mathbf{w}_I$. In practise we normalise $\mathbf{w}_S$ using (7) on each iteration.


The power iteration for the Geller-PageRank-Bollen-Eigenvector approach follows similarly. Define $\mathbf{D'} = \text{diag} (1/s_i)$ where $S_i$ is the row sum $s_i = \sum_k c_{ki}$ and the row-stochastic matrix $\Gamma' = \mathbf{K}\mathbf{D'}$. Then the PageRank power iteration is  
\begin{align}
	\mathbf{\pi^*}^{(k)} = \mathbf{M} \mathbf{\pi^*}^{(k-1)} \\
	\mathbf{M} = (1-\alpha)/N \mathbf{J} + \alpha \Gamma_{IS}^\top\Gamma_{SI}^T
\end{align}
where $\mathbf{J}$ is a matrix of all ones. The prestige is $\mathbf{w} = \mathbf{\pi^*}N\mathbf{D}$. The PageRank algorithm may be applied if the matrix $\Gamma_{SI}$ is not primitive.

\section{Application to Economics and Business}
We demonstrate and test JIRC using article and review data from the Web of Science (WoS) Core Collection in the ESI Field Economics \& Business (EB), published in the year 2020. The EB field encompasses four WoS categories: Business, Economics, Finance, and Management.

The dataset includes
\begin{itemize}
\item Journals (J): $M =613$ journals focused on Economics \& Business.
\item Institutions (I): $N =600$ institutions worldwide with a substantial number of publications in the ESI field Economics \& Business.
\item Published Articles (PA): More than $41,000$ articles published in 2020, classified within the Economics \& Business ESI Field by Clarivate's InCites and including journal names and institutional affiliations.
\item Citing Articles (CA): More than $200,000$ articles published between 2020 and 2025  classified within the Economics \& Business ESI Field by the Clarivate motor InCites, citing any of the  articles belonging to PA, including their full citation data and institutional affiliations.
\end{itemize}

Adopting institutional fractionation, we first construct $\mathbf{K_{IS}}$ as the ($N \times M$) matrix with elements $K_{IS}$($ij$) being the total number of citations from all articles in CA published by institution $i$ to all articles in PA published in journal $j$, and $\mathbf{K_{SI}}$ as the ($M \times N$) matrix with elements $K_{SI}$($ij$) being the total number of citations from all articles in CA published in journal $j$ to all articles published in PA by institution $i$. Following the PageRank methodology these matrices are made row-stochastic, creating transition matrices $\mathbf{\Gamma_{IS}^\top}$ and $\mathbf{\Gamma_{SI}^\top}$. We take $\alpha = 1$. 

To check the irreducibility of the two square matrices \(P_B^T P_A^T\) and \(P_A^T P_B^T\), we rely on a standard result in the field of matrix theory, connected to Perron-Frobenius theory, which states that a square matrix is irreducible if and only if its associated directed graph is strongly connected. We  used an algorithm based on \cite{Dijkstra} to settle the issue. Afterwards,
we solve for the principal eigenvector $v$ of the (N × N) matrix \(P_B^T P_A^T\) using the power method. The journal influence vector is then directly calculated as \(w = P_A^T v\).

\subsection{Researcher Screening}
To validate the practical utility of JIRC in individual research evaluation, we assembled a list of  330 prominent researchers in Economics \& Business, split into two groups:
\begin{enumerate}
  \item Control Group (C): $165$ researchers with universally acknowledged influence, including recipients of the Nobel Prize in Economics, the John Bates Clark Medal, Gossen Prize, and other prestigious field-specific awards, as well as European Research Council advanced grant holders.
  \item Test Group (T):$165$ researchers with a high number of highly-cited papers but without the prestigious awards defining the Control group.
\end{enumerate}
For each paper published by any of the researchers from our sample, we computed a JIRC score as the product of the journals's JIRC score and the best institutional JIRC score among all the institution in the paper byline

Then, for each researcher, we calculated a personal JIRC score by averaging the JIRC scores of all their publications between 2015 and 2024 (matching the  temporal window used by Clarivate in 2025 to identify highly cited papers). Researchers were then ranked and divided into three tiers based on this averaged score.

\section{Results}
\subsection{Journal and Institution Rankings}
The JIRC algorithm produced stable and convergent influence scores for all journals and institutions. The journal rankings ([Table/Appendix 1]) showed strong concordance with expert judgment and other reputation-based rankings, while also introducing nuance by factoring in the prestige of citing institutions. Similarly, the institutional rankings ([Table/Appendix 2]) reflected global academic reputations, effectively identifying centers of genuine intellectual influence.
\subsection{Disentangling Researcher's Influence from Notoriety}
The results of the researcher screening through individual JIRC scores were striking and demonstrate the core strength of the method.
\begin{itemize}
\item Tier 1 (Top 100 Researchers): This group was mainly composed  of researchers from the Control group $(88\%)$. These are individuals whose work is not only widely cited but is cited by influential entities (journals and institutions).
\item Tier 2 (Middle 130 Researchers): This group showed a mix of both Control and Test group members (75 vs 55), indicating a spectrum of influence among highly cited authors.
\item Tier 3 (Bottom 100 Researchers): This group was overwhelmingly populated by researchers from the Test group $(98\%)$. This suggests that while these individuals have achieved high citation counts (notoriety or popularity), their work is cited by sources with lower aggregate influence scores, potentially indicating strategic or circular citation practices within specific sub-communities, or work that is popular but not foundational.
    \end{itemize}
\section{Discussion}

Our proposed JIRC algorithm successfully extends the eigenvector centrality paradigm beyond the journal-journal citation network. By recirculating reputation through a journal-institution loop, we ground journal influence in the prestige of its citing institutions and, conversely, ground institutional prestige in the influence of its citing journals.

The most significant finding is the algorithm's ability to differentiate between influence and mere popularity at the individual researcher level. The clear separation between the award-winning Control group and the highly-cited-but-unawarded Test group across the tiers provides strong empirical validation. It suggests that JIRC is resistant to inflation from tactics like citation stacking or publishing in journals that prioritize rapid publication and high citation volume over intellectual rigor and influence.

The algorithm's accuracy is dependent on clean, disambiguated data for institutions. Future work will explore the impact of self-citations at the institutional and journal level, and the application of JIRC to other scientific fields and to author-level networks directly.
\section{Conclusion}

The Journal-Institution Reputation Circulation algorithm offers a robust and theoretically grounded framework for assessing scholarly influence. It moves beyond counting citations to weighing them based on a recursive definition of reputation that flows between producers (institutions) and disseminators (journals) of knowledge. Our results demonstrate that JIRC is not just a metric for ranking journals and institutions but also a powerful tool for research evaluation, capable of screening individual researchers to identify those with genuine, influential impact as opposed to those who have gained notoriety from a purely citation-based system. In an era of increasing publication volume and potential metric manipulation, JIRC provides a nuanced and reliable measure of true scholarly prestige.


\bibliographystyle{spbasic}
\bibliography{sciomtcs,MyLibrary}

\end{document}

